% !TEX root = ../elteikthesis_en.tex
\section{Implementation}
\label{sec:impl}

This section describes the full implementation of the ELTE IK Assistant across its four layers: data collection and ingestion, document embedding and vector storage, the RAG backend, and the web-based frontend. Each subsection covers one layer in detail, including the key design choices, data structures, and code.

% ─────────────────────────────────────────────────────────────
\subsection{Project Structure}
\label{sec:impl-structure}

The repository is organised into clearly separated concerns, as summarised in
Table~\ref{tab:structure}. The data pipeline scripts are never imported by the
serving layer; the \texttt{app/} package contains everything loaded at runtime.

\begin{longtable}{|l|p{0.58\textwidth}|}
    \hline
    \textbf{Path} & \textbf{Purpose} \\
    \hline \hline
    \endfirsthead
    \hline
    \textbf{Path} & \textbf{Purpose} \\
    \hline \hline
    \endhead
    \hline
    \endfoot

    \multicolumn{2}{|l|}{\textit{Setup and deployment}} \\
    \hline
    \texttt{setup.sh} / \texttt{setup.bat}  & First-run setup scripts for macOS/Linux and Windows. Pull Ollama models, create the virtual environment, run the crawler, and build the index. \\
    \hline
    \texttt{docker-compose.yml}             & Orchestrates the Ollama service, model pull, and FastAPI application as isolated containers for server deployment. \\
    \hline\hline
    \multicolumn{2}{|l|}{\textit{Data pipeline (run once; not imported at runtime)}} \\
    \hline
    \texttt{scripts/crawler.py}             & Resumable BFS web crawler for both ELTE domains. Downloads HTML, PDF, and DOCX files into \texttt{data/raw/}. \\
    \hline
    \texttt{scripts/build\_index.py}        & Production pipeline script. Ingests raw files, chunks them, generates embeddings with \texttt{all-MiniLM-L6-v2}, and upserts vectors into ChromaDB incrementally using the SHA-256 manifest. \\
    \hline\hline
    \multicolumn{2}{|l|}{\textit{Serving layer (\texttt{app/} package, loaded at runtime)}} \\
    \hline
    \texttt{app/settings.py}       & Centralised configuration constants: paths, model names, retrieval parameters, and timeouts. \\
    \hline
    \texttt{app/rag.py}            & Retrieval, prompt construction, and Ollama integration. Holds the ChromaDB client and embedding model as module-level singletons. \\
    \hline
    \texttt{app/ingest.py}         & File loading, chunking, and live embedding for user-uploaded documents. \\
    \hline
    \texttt{app/logger.py}         & SQLite-backed logging for chat interactions, sessions, and uploads. \\
    \hline
    \texttt{app/main.py}           & FastAPI application entry point. Wires all HTTP endpoints; contains no business logic. \\
    \hline
    \texttt{app/static/}           & Single-page HTML/CSS/JS frontend, served as a static file. \\
    \hline\hline
    \multicolumn{2}{|l|}{\textit{Development and evaluation (not part of the deployed system)}} \\
    \hline
    \texttt{scripts/chat\_cli.py}  & Interactive terminal client for the chat API. \\
    \hline
    \texttt{tests/}                & End-to-end pytest suite covering all API endpoints. \\
    \hline
    \texttt{notebooks/01--03}      & Exploration notebooks for ingestion, embedding, and RAG pipeline development. Superseded by \texttt{build\_index.py} for production use. \\
    \hline
    \texttt{notebooks/04\_evaluation.ipynb} & Evaluation of retrieval and answer quality across four model/RAG configurations. \\
    \hline\hline
    \multicolumn{2}{|l|}{\textit{Data directories (not tracked in version control)}} \\
    \hline
    \texttt{data/raw/}         & Crawled source documents (HTML; PDFs in \texttt{linked\_pdfs/}; DOCX in \texttt{linked\_docx/}). \\
    \hline
    \texttt{data/processed/}   & \texttt{chunks.json}, \texttt{manifest.json}, \texttt{pending\_changes.json}, and the ChromaDB vector store. \\
    \hline
    \texttt{data/logs/}        & SQLite chat log database. \\
    \hline
    \texttt{data/uploads/}     & User-uploaded documents stored after ingestion. \\
    \hline
    \caption{Project directory structure and file responsibilities.}
    \label{tab:structure}
\end{longtable}

% ─────────────────────────────────────────────────────────────
\subsection{Central Configuration}
\label{sec:impl-config}

All tunable parameters are centralised in \texttt{app/settings.py} so that
they can be changed in one place without touching any business logic. Every
other module imports from it; no constants are defined elsewhere.

\lstset{caption={Central configuration (\texttt{app/settings.py})},
        label=src:settings}
\begin{lstlisting}[language=Python]
from pathlib import Path

BASE_DIR             = Path(__file__).parent.parent
CHROMA_PATH          = str(BASE_DIR / "data/processed/chroma_db")
CHUNKS_PATH          = str(BASE_DIR / "data/processed/chunks.json")
MANIFEST_PATH        = str(BASE_DIR / "data/processed/manifest.json")
PENDING_CHANGES_PATH = str(BASE_DIR / "data/processed/pending_changes.json")
RAW_DIR              = str(BASE_DIR / "data/raw")
UPLOAD_DIR           = str(BASE_DIR / "data/uploads")
COLLECTION_NAME      = "elte_ik"
EMBEDDING_MODEL      = "all-MiniLM-L6-v2"
TOP_K                = 5
CHUNK_SIZE           = 800
CHUNK_OVERLAP        = 150
MAX_UPLOAD_BYTES     = 20 * 1024 * 1024   # 20 MB
ALLOWED_UPLOAD_TYPES = {
    "application/pdf", "text/html",
    "application/vnd.openxmlformats-officedocument.wordprocessingml.document",
    "application/msword", "application/octet-stream",
}
ALLOWED_EXTENSIONS   = {".pdf", ".html", ".htm", ".docx"}

OLLAMA_URL   = os.getenv("OLLAMA_URL", "http://localhost:11434/api/generate")
OLLAMA_MODEL = "llama3.2:3b"
TEMPERATURE  = 0.1
TIMEOUT_S    = 120
\end{lstlisting}

Table~\ref{tab:settings} groups the most commonly adjusted parameters by
concern.

\begin{table}[H]
    \centering
    \begin{tabular}{ | m{0.30\textwidth} | m{0.16\textwidth} | m{0.44\textwidth} | }
        \hline
        \textbf{Parameter} & \textbf{Default} & \textbf{Description} \\
        \hline \hline
        \multicolumn{3}{|l|}{\textit{Retrieval and chunking}} \\
        \hline
        \texttt{TOP\_K}        & \texttt{5}   & Chunks retrieved per query. \\
        \hline
        \texttt{CHUNK\_SIZE}   & \texttt{800} & Maximum chunk size in characters. \\
        \hline
        \texttt{CHUNK\_OVERLAP}& \texttt{150} & Overlap between adjacent chunks in characters. \\
        \hline
        \texttt{EMBEDDING\_MODEL} & \texttt{all-MiniLM-L6-v2} & Sentence embedding model. Must be identical between \texttt{build\_index.py} and the serving layer; a mismatch produces incorrect retrieval. \\
        \hline \hline
        \multicolumn{3}{|l|}{\textit{Language model}} \\
        \hline
        \texttt{OLLAMA\_MODEL} & \texttt{llama3.2:3b} & Default language model passed to Ollama. \\
        \hline
        \texttt{TEMPERATURE}   & \texttt{0.1} & LLM sampling temperature. \\
        \hline
        \texttt{TIMEOUT\_S}    & \texttt{120} & Per-request timeout for Ollama calls in seconds. \\
        \hline \hline
        \multicolumn{3}{|l|}{\textit{Upload validation}} \\
        \hline
        \texttt{MAX\_UPLOAD\_BYTES} & \texttt{20~MB} & Maximum accepted file size for user uploads. \\
        \hline
        \texttt{ALLOWED\_EXTENSIONS} & \texttt{.pdf .html .docx} & File extensions accepted by the upload endpoint. \\
        \hline
    \end{tabular}
    \caption{Key configuration parameters in \texttt{app/settings.py}.}
    \label{tab:settings}
\end{table}

\texttt{OLLAMA\_URL} is the only parameter read from the environment rather
than hardcoded. Its default value points to \texttt{localhost}, which is
correct for the manual setup path. The Docker Compose file overrides it with
\texttt{http://ollama:11434/api/generate} so that the application container
reaches the Ollama container by service name without any code change. This
single environment variable is what makes the same image work in both
deployment modes.

% ─────────────────────────────────────────────────────────────
\subsection{Data Collection and Ingestion}
\label{sec:impl-data}

ELTE Faculty of Informatics publishes its information across static HTML pages,
downloadable PDFs, and DOCX documents. The data pipeline collects and converts
all three formats into a uniform text representation before any machine learning
component is applied.

\subsubsection{File Loaders}

Each file type requires a different parsing strategy. PDFs are loaded with
PyMuPDF page by page. HTML files are parsed with \texttt{BeautifulSoup}:
structural noise tags (\texttt{script}, \texttt{style}, \texttt{nav},
\texttt{footer}) are stripped, then each paragraph is prefixed by its nearest
heading to preserve section context. DOCX files are parsed with
\texttt{python-docx}, extracting paragraphs and tables while tracking heading
structure. All formats share a cleaning step that collapses blank lines and
normalises whitespace.

\lstset{caption={Text cleaning and unified file loader
        (\texttt{scripts/build\_index.py})}, label=src:loader}
\begin{lstlisting}[language=Python]
def clean_text(text: str) -> str:
    text = re.sub(r"\n{3,}", "\n\n", text)
    text = re.sub(r"[ \t]+", " ", text)
    return text.strip()

def load_file(file_path: str | Path) -> List[Document]:
    suffix = Path(file_path).suffix.lower()
    if suffix == ".pdf":          return load_pdf_pymupdf(file_path)
    elif suffix in {".txt",".md"}: return load_txt(file_path)
    elif suffix in {".html",".htm"}: return load_html(file_path)
    elif suffix == ".docx":       return load_docx(file_path)
    else: raise ValueError(f"Unsupported file type: {suffix}")
\end{lstlisting}

Every document is represented as a LangChain \texttt{Document} carrying its
text and a metadata dictionary (file name, file type, page number, and
\texttt{source\_relative} path). This metadata is preserved into ChromaDB and
surfaced to the user as source references on every answer.

\subsubsection{Web Crawler}

Raw documents are collected by \texttt{scripts/crawler.py}, which performs a
resumable breadth-first crawl across two domains: all English pages under
\texttt{inf.elte.hu} recursively, and twenty whitelisted path subtrees on
\texttt{www.elte.hu} covering Neptun registration, visa procedures, student
finances, and dormitory housing. Explicit path prefixes on the second domain
prevent the crawler from wandering into unrelated university-wide content.

The crawler confirms each page is English via the \texttt{lang} attribute and
URL path before saving it. Linked PDFs and DOCX files are downloaded separately
into \texttt{linked\_pdfs/} and \texttt{linked\_docx/}. A 0.5-second delay
between requests limits server load. The crawler is resumable — existing files
are skipped — and a \texttt{--reset} flag starts fresh. The final corpus
contains 482 files.

\subsubsection*{Breadth-First Search Traversal}

BFS is implemented with a \texttt{deque} as the frontier queue. It explores
pages level by level, ensuring shallower and typically more important pages are
crawled before deep sub-pages, while bounding memory to the width of the graph
rather than its depth.

\begin{algorithm}[H]
\caption{BFS Web Crawler}
\label{alg:bfs-crawler}
\begin{algorithmic}[1]
\Require Seeds $S$, allowed-domain predicate $\textit{allowed}(\cdot)$
\State $\textit{visited} \gets \emptyset$
\State $\textit{queue} \gets \textsc{Deque}(S)$
\While{$\textit{queue}$ is not empty}
    \State $u \gets \textit{queue.popleft}()$
    \If{$u \in \textit{visited}$} \textbf{continue} \EndIf
    \State $\textit{visited} \mathrel{+}= \{u\}$
    \If{$u$ already on disk} parse cached HTML; \textbf{goto} link extraction \EndIf
    \State $\textit{html} \gets \textsc{Fetch}(u)$
    \If{$\textit{html}$ is English} save to \texttt{data/raw/} \EndIf
    \For{each link $v$ extracted from $\textit{html}$}
        \If{$v$ ends in \texttt{.pdf} or \texttt{.docx} \textbf{and} $\textit{allowed}(v)$}
            \State download $v$ to \texttt{linked\_pdfs/} or \texttt{linked\_docx/}
        \ElsIf{$\textit{allowed}(v)$ \textbf{and} $v \notin \textit{visited}$}
            \State $\textit{queue.append}(v)$
        \EndIf
    \EndFor
\EndWhile
\end{algorithmic}
\end{algorithm}

\subsubsection{Chunking Strategy}

A language model has a limited context window, so documents are split into
overlapping chunks using LangChain's \texttt{RecursiveCharacterTextSplitter}.

\lstset{caption={Chunking configuration (\texttt{scripts/build\_index.py})},
        label=src:chunking}
\begin{lstlisting}[language=Python]
text_splitter = RecursiveCharacterTextSplitter(
    chunk_size=800,
    chunk_overlap=150,
    length_function=len,
    separators=["\n\n", "\n", ". ", " ", ""]
)
\end{lstlisting}

800 characters is large enough to contain a complete regulation clause yet
small enough for the embedding model to encode with high fidelity. The
150-character overlap ensures that no sentence is silently dropped at a split
boundary. The separator list defines a priority order — paragraph breaks first,
then line breaks, then sentence ends — falling back to mid-word splits only as
a last resort. Chunks shorter than 100 characters are discarded, and each
surviving chunk is assigned a \texttt{chunk\_id} of the form
\texttt{<source\_relative>::chunk\_N}.

\subsubsection{Incremental Ingestion with SHA-256 Manifest}

Re-processing all 482 source files from scratch on every run would be wasteful.
\texttt{build\_index.py} therefore maintains a manifest at
\texttt{data/processed/manifest.json} recording the SHA-256 content hash,
chunk count, and timestamp for every ingested file.

On each run the script computes the hash of every file on disk and compares it
against the manifest, producing three change sets: new, updated, and deleted
files. Only new and updated files are re-chunked.

\lstset{caption={Incremental diff logic (\texttt{scripts/build\_index.py})},
        label=src:manifest}
\begin{lstlisting}[language=Python]
new_files     = sorted(disk_paths - manifest_paths)
deleted_files = sorted(manifest_paths - disk_paths)
updated_files = sorted(
    p for p in (disk_paths & manifest_paths)
    if manifest[p]["content_hash"] != current_files[p][1]
)
\end{lstlisting}

The diff results are written to \texttt{pending\_changes.json}. In Phase~2 of
the same run, the embedding step reads this file to determine exactly which
chunk IDs to delete from ChromaDB and which new vectors to insert, avoiding a
full re-embed of the collection. Once complete, \texttt{pending\_changes.json}
is deleted so the next run starts clean.
% ─────────────────────────────────────────────────────────────
\subsection{Embedding and Vector Storage}
\label{sec:impl-embedding}

Each text chunk is converted into a numeric vector — an
\emph{embedding} — that captures its semantic meaning. Chunks with similar
meaning produce similar vectors, allowing the retrieval step to find relevant
content even when the user's phrasing differs from the exact wording in the
source document.

\subsubsection{Embedding Model}

The chosen model is \texttt{all-MiniLM-L6-v2}, a compact sentence transformer
that produces 384-dimensional vectors and runs efficiently on CPU.
\texttt{build\_index.py} encodes all chunks in a single batched call:

\lstset{caption={Embedding pipeline (\texttt{scripts/build\_index.py})},
        label=src:embedding}
\begin{lstlisting}[language=Python]
class EmbeddingPipeline:
    def __init__(self, model_name: str = EMBEDDING_MODEL):
        self.model = SentenceTransformer(model_name)

    def encode(self, texts: list[str]) -> np.ndarray:
        return self.model.encode(texts, show_progress_bar=True)
\end{lstlisting}

The same model is loaded as a singleton in the serving layer so that document
and query embeddings always live in the same vector space. A mismatch between
the model used at index time and at query time produces silently incorrect
retrieval results.

\subsubsection{ChromaDB Vector Store}

The embeddings, texts, and metadata are stored in ChromaDB, a lightweight
vector database that persists to disk and requires no separate server process.
A single collection named \texttt{elte\_ik} is created with cosine similarity
as the distance metric.

\lstset{caption={Collection creation and upsert
        (\texttt{scripts/build\_index.py})}, label=src:chroma}
\begin{lstlisting}[language=Python]
client = chromadb.PersistentClient(path=CHROMA_PATH)
collection = client.get_or_create_collection(
    name=COLLECTION_NAME,
    metadata={"hnsw:space": "cosine"}
)
collection.upsert(
    ids=[str(c["metadata"]["chunk_id"]) for c in chunks],
    embeddings=embeddings.tolist(),
    documents=texts,
    metadatas=[c["metadata"] for c in chunks]
)
\end{lstlisting}

\texttt{upsert} rather than \texttt{add} ensures that re-running
\texttt{build\_index.py} after new documents have been crawled updates
existing chunks in place rather than creating duplicates.

\subsubsection*{Cosine Similarity}

For a query embedding $\mathbf{q}$ and a chunk embedding $\mathbf{d}$, cosine
similarity is defined as:

\begin{equation}
\text{sim}(\mathbf{q}, \mathbf{d}) =
    \frac{\mathbf{q} \cdot \mathbf{d}}{\|\mathbf{q}\|\,\|\mathbf{d}\|}
\label{eq:cosine}
\end{equation}

The result lies in $[-1, 1]$, where 1 means maximally similar and 0 means
unrelated. It is preferred over Euclidean distance for sentence embeddings
because it is invariant to vector magnitude — two chunks expressing the same
concept in different amounts of text score similarly, whereas Euclidean
distance would penalise the longer chunk for having a larger-norm embedding.

\subsubsection*{HNSW Index}

ChromaDB indexes the stored vectors using the Hierarchical Navigable Small
World (HNSW) graph~\cite{malkov2018hnsw}, which performs approximate
nearest-neighbour search in sub-linear time. HNSW builds a multi-layer
proximity graph and navigates it greedily from a random entry point,
converging on the nearest neighbours without examining every element. At
6128 stored chunks the brute-force alternative would also be fast enough, but
HNSW ensures the system remains performant as the collection grows.

% ─────────────────────────────────────────────────────────────
\subsection{Retrieval-Augmented Generation Pipeline}
\label{sec:impl-rag}

The RAG pipeline is the core intelligence of the system. It consists of three steps that run on every incoming user query: retrieval, prompt construction, and LLM generation.
All three are implemented in
\texttt{app/rag.py}.

\subsubsection{Retrieval}

The ChromaDB client, collection, and embedding model are initialised as
module-level singletons at import time, so they are loaded from disk once
when the FastAPI process starts and reused across all requests. Both
\texttt{SentenceTransformer} and ChromaDB's \texttt{PersistentClient} are
safe for concurrent read access, so no locking is required under
uvicorn's default threaded worker.

\lstset{caption={Singletons and retrieval function (\texttt{app/rag.py})},
        label=src:retrieve}
\begin{lstlisting}[language=Python]
_client     = chromadb.PersistentClient(path=CHROMA_PATH)
_collection = _client.get_collection(name=COLLECTION_NAME)
_model      = SentenceTransformer(EMBEDDING_MODEL)

def retrieve(query: str, n_results: int = TOP_K) -> list[dict]:
    query_emb = _model.encode([query]).tolist()
    results   = _collection.query(
        query_embeddings=query_emb, n_results=n_results
    )
    return [
        {"content": doc, "metadata": meta}
        for doc, meta in zip(
            results["documents"][0], results["metadatas"][0]
        )
    ]
\end{lstlisting}

When a query arrives, the embedding model encodes it into a 384-dimensional
vector and ChromaDB returns the \texttt{TOP\_K} (default: 5) chunks with the
highest cosine similarity to that vector.

\subsubsection{Prompt Construction}

The retrieved chunks are assembled into a structured prompt that instructs
the LLM to answer strictly from the provided context and to acknowledge when
the answer is absent. Each chunk is labelled with its source file name.

\lstset{caption={Prompt template (\texttt{app/rag.py})}, label=src:prompt}
\begin{lstlisting}[language=Python]
def build_prompt(query: str, chunks: list[dict]) -> str:
    context = "\n\n".join(
        f"[Source: {c['metadata']['file_name']}]\n{c['content']}"
        for c in chunks
    )
    return (
        "You are a helpful assistant for ELTE Faculty of Informatics "
        "students. Answer the question using only the context below. "
        "If the answer is not in the context, say so.\n\n"
        f"Context:\n{context}\n\n"
        f"Question: {query}\nAnswer:"
    )
\end{lstlisting}

The instruction \emph{``If the answer is not in the context, say so''} is
critical for a faculty assistant: acknowledging a gap in the indexed documents
is always preferable to hallucinating a regulation or deadline.

\subsubsection{LLM Generation via Ollama}

Answers are generated by a locally running model served through Ollama's HTTP
API. The default model is \texttt{llama3.2:3b}, but the system supports
runtime model switching — the selected model name is passed with each request
and forwarded to \texttt{call\_ollama}.

\lstset{caption={Ollama invocation (\texttt{app/rag.py})}, label=src:ollama}
\begin{lstlisting}[language=Python]
def call_ollama(prompt: str, model: str = OLLAMA_MODEL) -> str:
    if not check_ollama():
        raise RuntimeError("Ollama is not running.")
    payload = {
        "model": model, "prompt": prompt, "stream": False,
        "options": {"temperature": TEMPERATURE, "num_ctx": 2048},
    }
    try:
        r = requests.post(OLLAMA_URL, json=payload, timeout=TIMEOUT_S)
        r.raise_for_status()
        return r.json()["response"].strip()
    except requests.exceptions.Timeout:
        raise RuntimeError(f"Ollama timed out after {TIMEOUT_S}s")
    except requests.exceptions.HTTPError as e:
        raise RuntimeError(
            f"Ollama HTTP {e.response.status_code} -- "
            f"is model '{model}' pulled?"
        )
    except KeyError:
        raise RuntimeError(f"Unexpected Ollama response: {r.text[:200]}")
\end{lstlisting}

The context window is fixed at 2048 tokens, which fits five retrieved chunks
plus the prompt template with room for the generated answer. The
\texttt{try}/\texttt{except} block distinguishes three failure modes — timeout
(model still loading), HTTP error (model not pulled), and malformed response
body — each producing a descriptive \texttt{RuntimeError} that the API layer
converts into an HTTP~503.

\subsubsection{End-to-End Query Function}

The three steps are composed into \texttt{rag\_query}, the single entry point
called by the API layer:

\lstset{caption={End-to-end RAG query (\texttt{app/rag.py})},
        label=src:ragquery}
\begin{lstlisting}[language=Python]
def rag_query(query: str, model: str = OLLAMA_MODEL) -> dict:
    chunks = retrieve(query)
    prompt = build_prompt(query, chunks)
    answer = call_ollama(prompt, model=model)
    return {
        "answer": answer,
        "sources": [
            {"chunk_id": c["metadata"]["chunk_id"],
             "file":     c["metadata"]["file_name"]}
            for c in chunks
        ],
    }
\end{lstlisting}
% ─────────────────────────────────────────────────────────────
\subsection{FastAPI Backend}
\label{sec:impl-backend}

The backend is a FastAPI application (\texttt{app/main.py}) that exposes three HTTP endpoints and serves the static frontend files.

\begin{table}[H]
    \centering
    \begin{tabular}{|l|l|p{0.45\textwidth}|}
        \hline
        \textbf{Method} & \textbf{Path} & \textbf{Description} \\
        \hline \hline
        GET    & \texttt{/}                              & Serves the single-page chat interface (\texttt{index.html}). \\
        \hline
        GET    & \texttt{/health}                        & Returns application status and whether Ollama is reachable. \\
        \hline
        GET    & \texttt{/models}                        & Returns the list of models currently available in the local Ollama installation. \\
        \hline
        POST   & \texttt{/chat}                          & Accepts a user message and optional \texttt{session\_id} and \texttt{model}, runs the RAG pipeline, and returns the answer with source references. \\
        \hline
        GET    & \texttt{/sessions}                      & Returns the list of past chat sessions ordered by most recently updated. \\
        \hline
        GET    & \texttt{/sessions/\{id\}/messages}      & Returns all messages in a specific session in chronological order. \\
        \hline
        DELETE & \texttt{/sessions/\{id\}}               & Deletes a session and all its associated messages from the log. \\
        \hline
        POST   & \texttt{/upload}                        & Accepts a user-uploaded file, chunks and embeds it, and adds it to the live ChromaDB collection. \\
        \hline
    \end{tabular}
    \caption{API endpoints exposed by the FastAPI backend.}
    \label{tab:endpoints}
\end{table}

The \texttt{/chat} endpoint wraps \texttt{rag\_query} with error handling.

\lstset{caption={Chat endpoint (\texttt{app/main.py})}, label=src:chat-endpoint}
\begin{lstlisting}[language=Python]
@app.post("/chat", response_model=ChatResponse)
def chat(req: ChatRequest):
    try:
        result = rag_query(req.message)
    except RuntimeError as e:
        raise HTTPException(status_code=503, detail=str(e))
    return result
\end{lstlisting}

If \texttt{rag\_query} raises a \texttt{RuntimeError} (for example because Ollama is not running or timed out), the endpoint converts it into an HTTP 503 response so the frontend can display a meaningful error rather than a timeout spinner. CORS middleware is enabled to allow the frontend to communicate with the backend. For a locally deployed application this is straightforward, but in a production web deployment the allowed origins would need to be restricted to a specific domain.

% ─────────────────────────────────────────────────────────────
\subsection{Logging System}
\label{sec:impl-logging}

Every chat interaction is recorded in a SQLite database at \texttt{data/logs/chat\_logs.db}. The schema captures the timestamp, the user's message, the generated answer, the source references (as a JSON array), the response time in milliseconds, and any error message.

\lstset{caption={Database schema initialisation (\texttt{app/logger.py})}, label=src:schema}
\begin{lstlisting}[language=Python]
con.execute("""
    CREATE TABLE IF NOT EXISTS chat_logs (
        id            INTEGER PRIMARY KEY AUTOINCREMENT,
        timestamp     TEXT    NOT NULL,
        user_message  TEXT    NOT NULL,
        answer        TEXT    NOT NULL,
        sources       TEXT    NOT NULL,  -- JSON array
        response_ms   INTEGER NOT NULL,
        error         TEXT    DEFAULT NULL,
        session_id    TEXT    DEFAULT NULL
    )
""")
con.execute("""
    CREATE TABLE IF NOT EXISTS sessions (
        session_id  TEXT PRIMARY KEY,
        title       TEXT NOT NULL,
        created_at  TEXT NOT NULL,
        updated_at  TEXT NOT NULL
    )
""")
con.execute("""
    CREATE TABLE IF NOT EXISTS uploads (
        id          INTEGER PRIMARY KEY AUTOINCREMENT,
        timestamp   TEXT    NOT NULL,
        file_name   TEXT    NOT NULL,
        sha256      TEXT    NOT NULL,
        size_bytes  INTEGER NOT NULL,
        chunk_count INTEGER NOT NULL,
        status      TEXT    NOT NULL
    )
""")
\end{lstlisting}

The \texttt{chat\_logs} table links each message to a \texttt{session\_id}, enabling per-session history retrieval. The \texttt{sessions} table stores one row per conversation — its title (the first 60 characters of the opening message) and timestamps — so the sidebar can list sessions without scanning the full message log. The \texttt{uploads} table records every user-uploaded file alongside its SHA-256 digest and the number of chunks produced, providing an audit trail for the live document ingestion feature.


% ─────────────────────────────────────────────────────────────
\subsection{Frontend}
\label{sec:impl-frontend}

The user interface is a single-page HTML application served as a static file by FastAPI. It uses vanilla JavaScript without any framework dependency, which keeps the bundle small and removes the need for a separate build step.

The visual design follows the ELTE colour palette: navy blue (\texttt{\#012850}) for the header, teal (\texttt{\#009FA3}) for interactive elements, and light grey for the chat background. The layout consists of a fixed header bar, a scrollable message area, and a pinned input row at the bottom.

When the user sends a message, the frontend posts a JSON request to \texttt{/chat} and appends a typing indicator to the message area while waiting for the response. On success, both the answer text and the source file names are displayed in the assistant's reply bubble. On failure (HTTP 503), an error message is shown inline without leaving the page.

A status indicator in the header reflects the current Ollama availability by polling the \texttt{/health} endpoint on page load. This gives the user immediate feedback if the local model is not running before they attempt their first query.

% ─────────────────────────────────────────────────────────────

